\section{Critical defects and the emergence of the atom}
\label{rw:sec:defects}

\begin{theorem}[Completely monotone critical defects]\label{rw:thm:defect}
Let $a=1-b\in(0,1)$ and
\[
 D_n=\frac1a-\frac{H_{n-1}^{(b)}}{n^a},\qquad n\ge1.
\]
Then $D_n=\int_0^1u^{n-1}\dd\sigma_{a,b}(u)$ where $d\sigma_{a,b}(u)=-K_{1-b,b}(-\log u)\,du$ is the strictly positive critical negative part. If $\Delta D_n=D_{n+1}-D_n$,
then $(-1)^k\Delta^kD_n>0$ for every $k\ge0,n\ge1$. Moreover,
\begin{equation}\label{rw:eq:defectbounds}
 -\zeta(b)n^{-a}+\frac1{2n}<D_n
 <-\zeta(b)n^{-a}+\frac1n.
\end{equation}
Every finite Hankel matrix $(D_{j+k+1})_{j,k=0}^r$ is positive
definite, and $D_nD_{n+2}>D_{n+1}^2$.
\end{theorem}
\begin{proof}
Subtract the critical signed-measure moments from its atomic mass. Then
\[
 (-1)^k\Delta^kD_n=\int_0^1u^{n-1}(1-u)^k\dd\sigma(u)>0.
\]
On $w=1$, $k_w(T)=1$, so \eqref{rw:aux:Cbounds} reads
$1/2<C_{a,b}(T)<1$. Integrating these inequalities against
$e^{-nT}\dd T$ gives \eqref{rw:eq:defectbounds}.
For a nonzero real vector $(v_0,\ldots,v_r)$,
\[
 \sum_{j,k=0}^rv_jv_kD_{j+k+1}
 =\int_0^1\left(\sum_{j=0}^rv_ju^j\right)^2\dd\sigma(u)>0,
\]
since a nonzero polynomial cannot vanish on an interval and the
density is positive everywhere inside $(0,1)$. Strict log-convexity
follows either from its $2\times2$ minors or from strict
Cauchy--Schwarz with the weight $u^{n-1}\dd\sigma$.
\end{proof}

The theorem also gives a continuous version. For $q>0$,
\begin{equation}\label{rw:eq:continuousdefect}
 D(q)=\frac1a-q^{-a}\bigl(\zeta(b)-\zeta(b,q)\bigr)
 =\int_0^\infty e^{-qT}
   \bigl[-\zeta(b)k_a(T)+C_{a,b}(T)\bigr]\dd T.
\end{equation}
Thus $(-1)^jD^{(j)}(q)>0$ for all integers $j\ge0$.
All derivatives are justified by exponential domination. These are
inequalities for specific harmonic/Hurwitz expressions, not claims
about independence of their special values.

\section{Elementary critical-kernel identities and Hurwitz inequalities}
\label{rw:sec:criticalidentities}

In this section $0<b<1$, $a=1-b$, and
$C_b(T)=C_{1-b,b}(T)=\E[q(TV)]$ with
$V\sim\mathrm{Beta}(b,1-b)$. In particular $1/2<C_b(T)<1$ and
$C_b$ is strictly increasing.

\begin{theorem}[An elementary convergent series]\label{rw:thm:elementary}
For every $T>0$,
\begin{equation}\label{rw:eq:elementary}
 C_b(T)=\frac12+\sum_{k=1}^\infty
 \frac{\sin\!\bigl(b\arctan(T/(2\pi k))\bigr)}
 {\pi k\bigl(1+(T/(2\pi k))^2\bigr)^{b/2}}.
\end{equation}
Every summand is positive. The tail after $K\ge1$ terms is bounded by
\begin{equation}\label{rw:eq:elementarytail}
 0<C_b(T)-\frac12-\sum_{k=1}^K(\text{summand})
 \le\frac{bT}{2\pi^2K}.
\end{equation}
For $|T|<2\pi$ it also has the convergent expansion
\begin{equation}\label{rw:eq:Bernoulli}
 C_b(T)=\frac12+\sum_{j=1}^\infty
 \frac{B_{2j}(b)_{2j-1}}{(2j)!(2j-1)!}T^{2j-1},
\end{equation}
where $(b)_m$ is the rising factorial and $B_{2j}$ is a Bernoulli
number, not a Bernoulli polynomial evaluated at $b$.
\end{theorem}
\begin{proof}
The classical partial fraction for $\coth$ gives
\begin{equation}\label{rw:eq:hpartial}
 q(t)=\frac12+2t\sum_{k=1}^\infty\frac1{t^2+(2\pi k)^2}.
\end{equation}
This follows from DLMF 4.36.3 after rescaling
\cite{rw:DLMFHyperbolic}. The summands are positive for $t>0$, so
Tonelli permits averaging over $V$.
For $|z|<1$, the beta moments yield
\[
 \E\frac1{1-zV}=\sum_{j\ge0}z^j\frac{(b)_j}{j!}=(1-z)^{-b}.
\]
Both sides are analytic for $z\in\realcut$, so the same identity holds
there. At $z=\iu q$, its imaginary part is
\[
 \E\frac{qV}{1+q^2V^2}
 =\frac{\sin(b\arctan q)}{(1+q^2)^{b/2}}.
\]
Taking $q=T/(2\pi k)$ proves \eqref{rw:eq:elementary}.
Since $\sin(b\arctan q)\le bq$ and the denominator factor is at
least one, its tail is at most
$bT(2\pi^2)^{-1}\sum_{k>K}k^{-2}\le bT/(2\pi^2K)$.

The usual Bernoulli expansion of \eqref{subunit:eq:q} has radius $2\pi$:
$q(t)=1/2+\sum_{j\ge1}B_{2j}t^{2j-1}/(2j)!$.
For $|T|<2\pi$ it is uniformly absolutely convergent for
$0\le V\le1$. Average termwise and use
$\E V^{2j-1}=(b)_{2j-1}/(2j-1)!$ to obtain
\eqref{rw:eq:Bernoulli}.
\end{proof}

For instance,
\begin{equation}\label{rw:eq:Cexample}
 C_{1/2}(T)=\frac12+\frac{T}{24}-\frac{T^3}{2304}
                       +\frac{T^5}{122880}+O(T^7).
\end{equation}
The companion program regenerates exact coefficients through degree
$13$ at $b=1/10,1/2,9/10$ by two algebraically equivalent constructions.
The numerical tests of \eqref{rw:eq:elementary} use independent beta
quadrature, with the analytic tail bound displayed separately.

\begin{theorem}[A completely monotone normalized Hurwitz expression]
\label{rw:thm:hurwitzcm}
For $0<b<1$ and $q>0$, define
\begin{equation}\label{rw:eq:L}
 L_b(q)=q^{b-1}\zeta(b,q)+\frac1{1-b}.
\end{equation}
Then
\begin{equation}\label{rw:eq:criticalLaplace}
 L_b(q)=\int_0^\infty e^{-qT}C_b(T)\dd T,
 \qquad \frac1{2q}<L_b(q)<\frac1q.
\end{equation}
In particular, $(-1)^jL_b^{(j)}(q)>0$ for all $j\ge0$.
The function
\begin{equation}\label{rw:eq:normalizedHurwitz}
 qL_b(q)=q^b\zeta(b,q)+\frac q{1-b}
\end{equation}
decreases strictly from $1$ to $1/2$ as $q$ increases from zero to
infinity.
\end{theorem}
\begin{proof}
Formula \eqref{rw:eq:criticalLaplace} is \eqref{rw:aux:CLaplace} with
$a=1-b$. Bounds and strict complete monotonicity follow from
$1/2<C_b<1$ and differentiating the positive Laplace integral.
Substitution $T=s/q$ gives
\[
 qL_b(q)=\int_0^\infty e^{-s}C_b(s/q)\dd s.
\]
Strict increase of $C_b$ proves strict decrease with $q$.
The endpoint limits follow by dominated convergence and the limits
$C_b(0+)=1/2$, $C_b(\infty)=1$.
\end{proof}

The beta average is the common source of the elementary series and
the Hurwitz inequalities. The inequalities, finite-difference signs,
and Hankel determinants are analytic consequences of a positive
integral, not numerical pattern extrapolations.
